\section{Introduction}

Autonomous robotic systems are moving out of the laboratory and into shared public space~\cite{jung2018robots, pelikan2024encountering}. Delivery robots now roam city sidewalks~\cite{starship2025, yu2024understanding}, delivery drones ferry packages over residential neighborhoods~\cite{faa2025drone}, and driverless taxis carry passengers through downtown streets in cities around the world~\cite{waymo2025rides}. As these systems scale from isolated pilots toward routine presence, public frustration has also become visible. Residents have reported that delivery robots clutter street corners or block crosswalks~\cite{han2024codesign, yu2024understanding}, while some public encounters have involved obstruction or vandalism~\cite{sfstandard2024waymo, brscic2015escaping}. These encounters raise questions about how people read a robot's actions, work out what it is doing, and judge whether it belongs in a particular place.

Roboticists have long sought to deploy robots in public spaces, and HRI researchers have increasingly begun to examine how people might accommodate and make sense of these public robot encounters.~\cite{pelikan2024encountering,ekman2026sharing,bu2025making,weiss2010robots,jensen2005robots, thrun2000interaction} 
Pelikan et al. recently identified aspects of public sites and how they should inform robot design~\cite{pelikan2025localism}. However, it is often the case that a base robot platform is designed to be deployed in multiple locations. This work raises the question of whether the same robot remains equally legible across deployment sites. Comparing how people read the same service concept across settings can show which parts of an interpretation arise from the robot and which depend on the place in which it is encountered.

We pursue this question through a qualitative field study using two mobile Trashbots across five deployment sites spanning all five boroughs of New York City, building on prior field studies of mobile trash-receptacle robots~\cite{bu2024field}. We used a Wizard-of-Oz setup and conducted post-encounter interviews with people who engaged with the robots to examine how they read the robots' actions and made sense of their presence~\cite{bu2025making}. Our sites included a riverside park in Queens, a commercial corridor in the Bronx, a neighborhood park in Staten Island, a campus caf\'e in Manhattan, and an industrial and commercial campus in Brooklyn (Figure~\ref{fig:teaser}). Together, the sites varied in everyday use, sanitation arrangements, institutional context, and degree of public access. The Trashbots served as a common \emph{design probe}\cite{boehner201214} across these settings. The study took place across 22 field days, during which the Trashbots encountered approximately 490 people. We collected 86 post-encounter interview records involving 115 people, all of which were included in the comparative analysis (Table~\ref{tab:sites}).

Our study addresses three research questions:


\begin{itemize}

\item \textbf{RQ1:} How do people ``read'' public service robots across our different sites?

\item \textbf{RQ2:} What strategies do people use to make sense of the robots' presence in their space?

\item \textbf{RQ3:} What factors influence how people read the robots' actions? Are they specific to the encounter site?

\end{itemize}

Our qualitative analysis examined how participants interpreted the Trashbots, the strategies they used to make sense of the robots' presence, and what informed these interpretations. We also drew on available onboard-camera footage to  characterize how people interacted with the robots during encounters.


Our findings motivate the development of the concept of \emph{local legibility} as a design objective for public service robots. We use the term to describe whether people can read a robot's actions and understand its role in relation to the place where they encounter it. This includes understanding what service the robot provides, who is responsible for it, what its sensing practices mean, and how it relates to existing activities, infrastructure, and services at the site. 

This work makes these contributions: \begin{itemize}  

\item We provide field observations of public robot encounters and post-encounter interviews
    across five sites in New York City. 
    
    \item We characterize the interpretive process through which people
make sense of public robots, identifying the shared questions
that structure this process across deployment settings. 

\item We show the knowledge and resource people draw on to address these interpretive process, producing different judgments about the robots across deployment settings.

 \item We derive design implications for local legibility: making robots’ purpose, institutional responsibility, and recording practices understandable within their settings, and adapting behavior and siting to local activities and needs.  \end{itemize}